% MAT670 -- Interlude after Lecture 6: from Dynkin diagrams to Lie algebras and Lie groups.
% Entirely new material (not in the handwritten notes). Labels prefixed LIE:.

\chapter[Interlude: Lie algebras and Lie groups]{Interlude: From Dynkin diagrams to Lie algebras and Lie groups}
\label{LIE:ch}

\begin{addedblock}
This interlude is new: it is not in the handwritten notes. It picks up the remark at the start of
Lecture~4, that the classification of root systems is ``half of the classification of complex
semisimple Lie algebras'', and follows it. To keep it readable it is set in black; all of it should be
read as added material.
\end{addedblock}

In Lectures~4--6 we took the root system axioms as given, classified the possible Dynkin diagrams, and
then built each root system by hand from lattices. Two questions were left hanging.
\begin{itemize}
\item Where do the axioms come from? Why should anyone care about a finite set of vectors closed under
  reflections, with the integrality condition $n_{\beta\alpha}\in\Z$?
\item The same root system $A_{n-1}$ appears twice in these notes, once as the vectors
  $e_i-e_j$ in the hyperplane $\sum x_i=0$ (Lecture~6) and, later, once as the walls $x_i=x_j$ of
  $\Conf(N,\R)$ (Lecture~8). Is there a single object behind both?
\end{itemize}
The answer to both is \emph{matrices}. In this interlude we find root systems inside matrix
algebras, see that the crystallographic axioms are forced by a small piece of linear algebra, and
then use the exponential map to pass from the algebra to groups of matrices. At the end the root
lattice comes back, this time sitting inside a group.

Throughout, $M_n$ denotes the $n\times n$ complex matrices and $E_{ij}$ the matrix with a $1$ in
position $(i,j)$ and zeros elsewhere, so that
\[
  E_{ij}E_{kl}=\delta_{jk}E_{il}.
\]

%======================================================================
\section{The roots of \texorpdfstring{$A_{n-1}$}{A(n-1)} inside matrices}
\label{LIE:sec:sln}

Let $h=\operatorname{diag}(h_1,\dots,h_n)$ be a diagonal matrix. Multiplying a matrix unit by $h$ on
the left scales its row and on the right scales its column, so
\[
  hE_{ij}=h_iE_{ij},\qquad E_{ij}h=h_jE_{ij},
  \qquad\text{hence}\qquad
  hE_{ij}-E_{ij}h=(h_i-h_j)\,E_{ij}.
\]
In words: the linear map
\[
  \operatorname{ad}_h\colon M_n\to M_n,\qquad X\mapsto [h,X]:=hX-Xh
\]
is \emph{diagonal} in the basis $\{E_{ij}\}$, with eigenvalue $h_i-h_j$ on $E_{ij}$.

Now regard the eigenvalue as a function of $h$. On $E_{ij}$ it is the linear functional
$h\mapsto h_i-h_j$. If we identify functionals with vectors using the standard inner product, so that
$h\mapsto h_i$ corresponds to $e_i$, then the eigenvalue on $E_{ij}$ is the vector $e_i-e_j$. So:
\begin{center}
\emph{the nonzero eigenvalues of $\operatorname{ad}_h$, as $h$ runs over diagonal matrices, are exactly
the roots $e_i-e_j$ of $A_{n-1}$.}
\end{center}
The diagonal matrices themselves have eigenvalue $0$.

Where is the hyperplane $\sum x_i=0$? Scalar matrices $h=cI$ commute with everything, so they tell us
nothing, and it is natural to discard them. That leaves the \emph{traceless} matrices
\[
  \mathfrak{sl}_n=\{X\in M_n:\operatorname{tr}X=0\},
\]
whose diagonal part
\[
  \mathfrak h=\{\operatorname{diag}(h_1,\dots,h_n): h_1+\dots+h_n=0\}
\]
is exactly the hyperplane $\sum x_i=0$ of Lecture~6, of dimension $n-1$. We get a decomposition into
eigenspaces
\begin{equation}\label{LIE:eq:rootdecomp}
  \mathfrak{sl}_n=\mathfrak h\ \oplus\ \bigoplus_{\alpha\in\Delta}\mathfrak g_\alpha,
  \qquad
  \mathfrak g_{e_i-e_j}=\C E_{ij},
\end{equation}
where $\Delta=\{e_i-e_j: i\neq j\}$ is the root system $A_{n-1}$. As a check on dimensions:
$n^2-1=(n-1)+n(n-1)$, which is (rank) $+$ (number of roots), using the count
$\abs{\Delta}=n(n-1)$ from Lecture~6.

\begin{example}[$\mathfrak{sl}_3$]\label{LIE:ex:sl3}
Here $\mathfrak h$ is two-dimensional and the six off-diagonal matrix units sit at the six roots of
$A_2$. Placing each $E_{ij}$ at its root gives the familiar hexagon:
\begin{center}
\begin{tikzpicture}[scale=1.55]
  \foreach \a/\lab/\pos in {0/E_{12}/right, 60/E_{13}/above right, 120/E_{23}/above left,
                            180/E_{21}/left, 240/E_{31}/below left, 300/E_{32}/below right}
    {\draw[-{Stealth}] (0,0)--(\a:1); \fill (\a:1) circle (0.03);
     \node[\pos,font=\small] at (\a:1) {$\lab$};}
  \fill (0,0) circle (0.045); \fill (0.09,0) circle (0.045);
  \node[below,font=\small] at (0.05,-0.08) {$\mathfrak h$};
  \draw[dashed,-{Stealth},addcol,thick] (120:1.05) to[bend left=35] node[above,font=\scriptsize]{$\operatorname{ad}_{E_{12}}$} (60:1.05);
  \node[font=\scriptsize,gray] at (2.6,0.55) {$\alpha_1=e_1-e_2$};
  \node[font=\scriptsize,gray] at (2.6,0.25) {$\alpha_2=e_2-e_3$};
  \node[font=\scriptsize,gray] at (2.6,-0.05) {$\alpha_1+\alpha_2=e_1-e_3$};
\end{tikzpicture}
\end{center}
The two dots at the centre stand for the two-dimensional zero eigenspace $\mathfrak h$. The dashed
arrow anticipates Section~\ref{LIE:sec:bracket}: bracketing with $E_{12}$ moves $E_{23}$ to $E_{13}$,
that is, it \emph{adds the root} $\alpha_1$.
\end{example}

The choices made in Lecture~4 also have matrix meanings. Take the direction $d=(n,n-1,\dots,1)$, so
that $\ip{e_i-e_j}{d}=j-i$. Then:
\begin{itemize}
\item the positive roots $e_i-e_j$ with $i<j$ are the matrix units \emph{above} the diagonal, and
  the negative roots are those below it;
\item the simple roots $\alpha_i=e_i-e_{i+1}$ are the entries just above the diagonal,
  $E_{12},E_{23},\dots,E_{n-1,n}$;
\item every positive root is a sum of simple roots: $e_i-e_j=\alpha_i+\alpha_{i+1}+\dots+\alpha_{j-1}$.
  On the matrix side this is the statement that $E_{ij}$ is an iterated product (indeed an iterated
  commutator) of superdiagonal units.
\end{itemize}
The Weyl group, generated by the reflections $s_{e_i-e_j}$, acts on $\mathfrak h$ by swapping $h_i$ and
$h_j$. It is the symmetric group $\Sigma_n$ permuting the diagonal entries, which is conjugation by
permutation matrices. The Weyl chamber containing $d$ is $\{h_1>h_2>\dots>h_n\}$, one of the $n!$
chambers of $\Conf(n,\R)$ from Lecture~8.

%======================================================================
\section{The bracket and the definition of a Lie algebra}
\label{LIE:sec:bracket}

What made the computation work was the operation $[X,Y]=XY-YX$, the \emph{commutator}. It has three
properties that are easy to check:
\begin{enumerate}[label=(\roman*)]
\item it is bilinear;
\item it is antisymmetric, $[X,Y]=-[Y,X]$;
\item it satisfies the \emph{Jacobi identity}
  \[
    [X,[Y,Z]]=[[X,Y],Z]+[Y,[X,Z]],
  \]
  which says that $\operatorname{ad}_X=[X,\cdot\,]$ obeys the product rule (it is a derivation of the
  bracket). Expanding both sides into products of three matrices proves it.
\end{enumerate}

\begin{definition}\label{LIE:def:liealg}
A \emph{Lie algebra} over a field $k$ is a vector space $\mathfrak g$ with a bilinear map
$[\cdot,\cdot]\colon\mathfrak g\times\mathfrak g\to\mathfrak g$ that is antisymmetric and satisfies the
Jacobi identity. A \emph{subalgebra} is a subspace closed under the bracket; an \emph{ideal} is a
subspace $I$ with $[\mathfrak g,I]\subseteq I$. A Lie algebra is \emph{simple} if it is not abelian and
has no ideals other than $0$ and $\mathfrak g$. A \emph{representation} of $\mathfrak g$ on a vector space
$V$ is a linear map $\rho\colon\mathfrak g\to\operatorname{End}(V)$ with
$\rho([X,Y])=\rho(X)\rho(Y)-\rho(Y)\rho(X)$.
\end{definition}

\begin{example}\leavevmode
\begin{enumerate}[nosep]
\item $\mathfrak{gl}_n=M_n$ with the commutator.
\item $\mathfrak{sl}_n$ is a subalgebra, since $\operatorname{tr}(XY-YX)=0$ for all $X,Y$. It is simple for
  $n\ge2$.
\item $\mathfrak{so}_n=\{X\in M_n: X^T=-X\}$ is a subalgebra: if $X,Y$ are antisymmetric, then
  $[X,Y]^T=Y^TX^T-X^TY^T=YX-XY=-[X,Y]$.
\item $\R^3$ with the cross product is a Lie algebra (Jacobi is the identity
  $a\times(b\times c)=(a\times b)\times c+b\times(a\times c)$). Sending $\omega\in\R^3$ to the matrix
  $X_\omega$ with $X_\omega v=\omega\times v$ gives an isomorphism $(\R^3,\times)\cong\mathfrak{so}_3(\R)$.
\item By the Jacobi identity, $X\mapsto\operatorname{ad}_X$ is a representation of $\mathfrak g$ on itself,
  the \emph{adjoint representation}. Section~\ref{LIE:sec:sln} was a study of the adjoint
  representation of $\mathfrak{sl}_n$, restricted to $\mathfrak h$.
\end{enumerate}
\end{example}

The Jacobi identity also explains the dashed arrow in Example~\ref{LIE:ex:sl3}. Suppose
$[h,x]=\alpha(h)x$ and $[h,y]=\beta(h)y$ for all $h\in\mathfrak h$. Then
\[
  [h,[x,y]]=[[h,x],y]+[x,[h,y]]=\big(\alpha(h)+\beta(h)\big)[x,y].
\]

\begin{lemma}\label{LIE:lem:add}
$[\mathfrak g_\alpha,\mathfrak g_\beta]\subseteq\mathfrak g_{\alpha+\beta}$, where
$\mathfrak g_{\alpha+\beta}=0$ if $\alpha+\beta$ is neither a root nor $0$, and $\mathfrak g_0=\mathfrak h$.
\end{lemma}

So the bracket moves around the root diagram by \emph{adding roots}. In $\mathfrak{sl}_3$,
$[E_{12},E_{23}]=E_{13}$ is the statement $\alpha_1+\alpha_2=e_1-e_3$, while $[E_{12},E_{13}]=0$
because $2e_1-e_2-e_3$ is not a root.

%======================================================================
\section{The atom: \texorpdfstring{$\mathfrak{sl}_2$}{sl(2)}}
\label{LIE:sec:sl2}

Fix a root $\alpha=e_i-e_j$ of $A_{n-1}$ and put
\[
  e_\alpha=E_{ij},\qquad f_\alpha=E_{ji},\qquad h_\alpha=[e_\alpha,f_\alpha]=E_{ii}-E_{jj}\in\mathfrak h.
\]
A short computation gives
\begin{equation}\label{LIE:eq:sl2rel}
  [h_\alpha,e_\alpha]=2e_\alpha,\qquad [h_\alpha,f_\alpha]=-2f_\alpha,\qquad [e_\alpha,f_\alpha]=h_\alpha .
\end{equation}
For $n=2$ these three matrices
\[
  e=\begin{pmatrix}0&1\\0&0\end{pmatrix},\quad
  f=\begin{pmatrix}0&0\\1&0\end{pmatrix},\quad
  h=\begin{pmatrix}1&0\\0&-1\end{pmatrix}
\]
are a basis of $\mathfrak{sl}_2$. So every root $\alpha$ comes with a copy
$\mathfrak s_\alpha=\operatorname{span}(e_\alpha,f_\alpha,h_\alpha)\cong\mathfrak{sl}_2$ sitting inside $\mathfrak{sl}_n$.

Now look at how $h_\alpha$ acts on the other root spaces. For a root $\beta=e_k-e_l$,
\[
  \beta(h_\alpha)=(\delta_{ki}-\delta_{kj})-(\delta_{li}-\delta_{lj})=\ip{\beta}{\alpha}
  =\frac{2\ip{\beta}{\alpha}}{\ip{\alpha}{\alpha}}=n_{\beta\alpha},
\]
since $\ip{\alpha}{\alpha}=2$. So \emph{the Cartan integer $n_{\beta\alpha}$ is the eigenvalue of
$h_\alpha$ on $\mathfrak g_\beta$}. ($h_\alpha$ is called the \emph{coroot} of $\alpha$.) To see why these
eigenvalues have to be integers, we need one theorem about $\mathfrak{sl}_2$.

\begin{theorem}[Representations of $\mathfrak{sl}_2$]\label{LIE:thm:sl2}
Let $V$ be a finite-dimensional complex representation of $\mathfrak{sl}_2$. Then $h$ acts
diagonalizably on $V$, its eigenvalues are integers, and for every $k\in\Z$ the eigenspaces satisfy
$\dim V_k=\dim V_{-k}$. If $V$ is irreducible of dimension $m+1$, then the eigenvalues of $h$ are
$m,m-2,\dots,-m$, each with multiplicity one.
\end{theorem}

\begin{proof}[Proof for irreducible $V$]
The key observation is that $e$ raises $h$-eigenvalues by $2$ and $f$ lowers them by $2$. If
$hv=\lambda v$, then by \eqref{LIE:eq:sl2rel}
\[
  h(ev)=e(hv)+[h,e]v=(\lambda+2)\,ev,\qquad h(fv)=(\lambda-2)\,fv .
\]
Take any eigenvector of $h$ and apply $e$ repeatedly. The vectors obtained are eigenvectors for
distinct eigenvalues while they are nonzero, and $V$ is finite-dimensional, so eventually we reach a
\emph{highest weight vector}: $v\neq0$ with $hv=\lambda v$ and $ev=0$. Put $v_k=f^kv$. Then
$hv_k=(\lambda-2k)v_k$, and induction on $k$ using $ef=fe+h$ gives
\[
  e\,v_k=k(\lambda-k+1)\,v_{k-1}.
\]
Let $m$ be the largest index with $v_m\neq0$. Then $0=e\,v_{m+1}=(m+1)(\lambda-m)\,v_m$, so
$\lambda=m$. In particular $\lambda$ is a non-negative integer. The span of $v_0,\dots,v_m$ is closed
under $e$, $f$ and $h$, so by irreducibility it is all of $V$. The eigenvalues are therefore
$m,m-2,\dots,-m$, as claimed.
\end{proof}

\begin{figure}[ht]
\centering
\begin{tikzpicture}[xscale=1.7]
  \foreach \x/\l in {0/{-m},1/{-m+2},2/{\cdots},3/{m-2},4/{m}} {\fill (\x,0) circle (0.05); \node[font=\small] at (\x,-0.62) {$\l$};}
  \foreach \x in {0,1,2,3} {\draw[-{Stealth},addcol] (\x+0.1,0.1) to[bend left=40] (\x+0.9,0.1);
                            \draw[-{Stealth},flagcol] (\x+0.9,-0.1) to[bend left=40] (\x+0.1,-0.1);}
  \node[addcol,font=\small] at (2,0.7) {$e$ raises by $2$};
  \node[flagcol,font=\small] at (2,-1.05) {$f$ lowers by $2$};
\end{tikzpicture}
\caption{The $(m+1)$-dimensional irreducible representation of $\mathfrak{sl}_2$: a string of
$h$-eigenvalues, symmetric about $0$.}
\end{figure}

For a general finite-dimensional $V$ one uses \emph{Weyl's complete reducibility theorem}: $V$ is a
direct sum of irreducible representations (see Humphreys, \emph{Introduction to Lie Algebras and
Representation Theory}, \S6.3, or Hall, \emph{Lie Groups, Lie Algebras, and Representations}, Chs.~4 and~10).
The statements about eigenvalues then follow from the irreducible case.

%======================================================================
\section{Why the root system axioms hold}
\label{LIE:sec:axioms}

Now apply Theorem~\ref{LIE:thm:sl2} to the copy $\mathfrak s_\alpha\cong\mathfrak{sl}_2$, acting by
brackets. Fix roots $\beta\neq\pm\alpha$ and form the \emph{$\alpha$-string through $\beta$},
\[
  V=\bigoplus_{k\in\Z}\mathfrak g_{\beta+k\alpha}.
\]
By Lemma~\ref{LIE:lem:add}, bracketing with $e_\alpha$ or $f_\alpha$ shifts $k$ by $\pm1$, so $V$ is a
representation of $\mathfrak s_\alpha$. On the summand $\mathfrak g_{\beta+k\alpha}$, $h_\alpha$ acts by the
scalar
\[
  (\beta+k\alpha)(h_\alpha)=n_{\beta\alpha}+2k .
\]

\begin{itemize}
\item \textbf{Axiom 4 (crystallographic).} $n_{\beta\alpha}$ is an eigenvalue of $h_\alpha$ on $V$
  (take $k=0$), so by Theorem~\ref{LIE:thm:sl2} it is an integer.
\item \textbf{Axiom 3 (reflections).} The eigenvalues are symmetric about $0$, so $-n_{\beta\alpha}$ is
  also an eigenvalue: $n_{\beta\alpha}+2k=-n_{\beta\alpha}$ for some $k$ with
  $\mathfrak g_{\beta+k\alpha}\neq0$. That is, $k=-n_{\beta\alpha}$, and
  \[
    \beta-n_{\beta\alpha}\,\alpha=s_\alpha(\beta)
  \]
  is a root. The reflection $s_\alpha$ maps $\Delta$ to itself.
\item \textbf{Strings have no gaps.} Each irreducible piece of $V$ has eigenvalues forming an unbroken
  string with step $2$. In fact $V$ is irreducible here, since the root spaces turn out to be
  one-dimensional. So the roots $\beta+k\alpha$ form an unbroken string $\beta-p\alpha,\dots,\beta+q\alpha$
  with $p-q=n_{\beta\alpha}$. This is the ``$\alpha-\beta$'' phenomenon behind Lemma~\ref{L04a:lem:diff} of Lecture~4.
\end{itemize}

For $\mathfrak{sl}_n$ all of this can of course be checked directly. The point is that the argument used
only three ingredients: a commutative subalgebra $\mathfrak h$ acting diagonalizably, the copies of
$\mathfrak{sl}_2$ attached to the roots, and finite-dimensionality. These are exactly what a
\emph{complex semisimple Lie algebra} provides. In such an algebra a maximal commutative subalgebra of
diagonalizable elements (a \emph{Cartan subalgebra}) always exists, and its nonzero eigenvalues form a
reduced, crystallographic root system. (Reducedness, axiom~2, takes one more argument; see Humphreys
\S8.4.) \emph{The root system axioms are the shadow of $\mathfrak{sl}_2$ representation theory.}

%======================================================================
\section{The other classical families}
\label{LIE:sec:classical}

The families $B_n$, $C_n$, $D_n$ from Lecture~6 come from the Lie algebras that preserve a bilinear
form. Let $J$ be an invertible $N\times N$ matrix and put
\[
  \mathfrak g_J=\{X\in M_N: X^TJ+JX=0\}.
\]
This is a subalgebra (the same computation as for $\mathfrak{so}_n$), and in Section~\ref{LIE:sec:exp}
we will see that it consists of the \emph{infinitesimal symmetries} of the form $\ip{v}{w}_J=v^TJw$.
For $J=I$ this is $\mathfrak{so}_N$.

The trick is to choose $J$ so that a large family of \emph{diagonal} matrices lies in $\mathfrak g_J$.
Write $i'=N+1-i$ and let $J$ be the antidiagonal matrix with $J_{i,i'}=1$. For $X\in M_N$ one computes
$(JX)_{ij}=X_{i'j}$ and $(X^TJ)_{ij}=X_{j'i}$, so
\[
  X\in\mathfrak g_J\iff X_{ij}=-X_{j'i'}\ \text{for all } i,j .
\]
(Over $\C$ any two nondegenerate symmetric forms are equivalent, so $\mathfrak g_J\cong\mathfrak{so}_N(\C)$.)
\begin{itemize}
\item The diagonal matrices in $\mathfrak g_J$ are
  $h=\operatorname{diag}(h_1,\dots,h_n,[0],-h_n,\dots,-h_1)$, where the middle $0$ is present only when
  $N=2n+1$ is odd. These form a commutative subalgebra $\mathfrak h$ of dimension $n$.
\item The matrices $E_{ij}-E_{j'i'}$, for pairs with $i\neq j$ and $j\neq i'$, span the rest of
  $\mathfrak g_J$ (for $j=i'$ the expression is $0$: the antidiagonal of $X$ vanishes). Each is
  an eigenvector of $\operatorname{ad}_h$ with eigenvalue $\tilde h_i-\tilde h_j$, where $\tilde h$ is the
  diagonal of $h$. (The two matrix units give the same eigenvalue, because $\tilde h_{k'}=-\tilde h_k$.)
\item Reading off the eigenvalues: $i,j\le n$ gives $\pm(e_i-e_j)$; $i\le n<j$ gives
  $\pm(e_i+e_k)$ with $j=k'$ and $k\neq i$ (the case $k=i$ is excluded since then $j=i'$); and when $N$
  is odd the middle index gives $\pm e_i$.
\end{itemize}
So $\mathfrak{so}_{2n}$ has root system $\{\pm e_i\pm e_j\}=D_n$ and $\mathfrak{so}_{2n+1}$ has
$\{\pm e_i\pm e_j\}\cup\{\pm e_i\}=B_n$. If instead
$J=\left(\begin{smallmatrix}0&K\\-K&0\end{smallmatrix}\right)$, with $K$ the $n\times n$ antidiagonal
matrix of ones, the form is \emph{symplectic}. Now signs enter: with $\varepsilon_i=1$ for $i\le n$
and $\varepsilon_i=-1$ for $i>n$, the condition becomes $X_{ij}=-\varepsilon_i\varepsilon_jX_{j'i'}$, the
diagonal part is the same $\mathfrak h$ as before, and the off-diagonal basis is
$E_{ij}-\varepsilon_i\varepsilon_jE_{j'i'}$. The same bookkeeping gives the symplectic algebra
$\mathfrak{sp}_{2n}$, whose root system is $\{\pm e_i\pm e_j\}\cup\{\pm 2e_i\}=C_n$: now the case $k=i$ is
allowed, since for $j=i'$ the basis element is $2E_{ii'}\neq0$, and it produces the long roots $\pm2e_i$.

\begin{center}
\renewcommand{\arraystretch}{1.25}
\begin{tabular}{llllll}
\toprule
type & Lie algebra & condition & roots & rank $+\abs{\Delta}$ & $\dim\mathfrak g$\\
\midrule
$A_n$ & $\mathfrak{sl}_{n+1}$ & $\operatorname{tr}X=0$ & $e_i-e_j$ & $n+n(n+1)$ & $(n+1)^2-1$\\
$B_n$ & $\mathfrak{so}_{2n+1}$ & $X^TJ+JX=0$ & $\pm e_i\pm e_j,\ \pm e_i$ & $n+2n^2$ & $n(2n+1)$\\
$C_n$ & $\mathfrak{sp}_{2n}$ & $X^TJ+JX=0$ ($J$ skew) & $\pm e_i\pm e_j,\ \pm2e_i$ & $n+2n^2$ & $n(2n+1)$\\
$D_n$ & $\mathfrak{so}_{2n}$ & $X^TJ+JX=0$ & $\pm e_i\pm e_j$ & $n+2n(n-1)$ & $n(2n-1)$\\
\bottomrule
\end{tabular}
\end{center}
In every row, $\dim\mathfrak g=\operatorname{rank}+\abs{\Delta}$, as it must be by the root decomposition
\eqref{LIE:eq:rootdecomp}. The root counts are the ones computed in Lecture~6.

\begin{remark}[The ``repetitions'' in the classification]
Theorem~\ref{L04a:thm:classification} restricted the ranges of $n$ ($B_n$ for $n\ge2$, $C_n$ for $n\ge3$, $D_n$ for $n\ge4$) to
avoid repeating diagrams. Each coincidence of small diagrams is an isomorphism of Lie algebras:
$B_1=C_1=A_1$ gives $\mathfrak{so}_3\cong\mathfrak{sp}_2=\mathfrak{sl}_2$; $B_2=C_2$ gives
$\mathfrak{so}_5\cong\mathfrak{sp}_4$; $D_2=A_1\times A_1$ gives $\mathfrak{so}_4\cong\mathfrak{sl}_2\times\mathfrak{sl}_2$; and
$D_3=A_3$ gives $\mathfrak{so}_6\cong\mathfrak{sl}_4$.
\end{remark}

%======================================================================
\section{The classification of simple Lie algebras}
\label{LIE:sec:classification}

Everything so far points to a dictionary between Dynkin diagrams and Lie algebras. The result is
due to Killing (1888--90) and \'E.~Cartan (thesis, 1894); the diagrams themselves are due to
Dynkin (1946).

\begin{theorem}[Killing--Cartan classification]\label{LIE:thm:KC}
Let $\mathfrak g$ be a complex semisimple Lie algebra.
\begin{enumerate}[nosep]
\item $\mathfrak g$ has a Cartan subalgebra $\mathfrak h$, unique up to automorphisms of $\mathfrak g$, and the
  nonzero eigenvalues of $\operatorname{ad}\mathfrak h$ form a root system $\Delta(\mathfrak g)$ in (the real
  span of) $\mathfrak h^*$, with one-dimensional root spaces.
\item $\mathfrak g$ is simple if and only if $\Delta(\mathfrak g)$ is irreducible.
\item $\mathfrak g\cong\mathfrak g'$ if and only if $\Delta(\mathfrak g)\cong\Delta(\mathfrak g')$.
\item (Existence; Serre's theorem) Every root system arises this way.
\end{enumerate}
Consequently the complex simple Lie algebras correspond exactly to the connected Dynkin diagrams
$A_n$ ($n\ge1$), $B_n$ ($n\ge2$), $C_n$ ($n\ge3$), $D_n$ ($n\ge4$), $E_6$, $E_7$, $E_8$, $F_4$, $G_2$.
\end{theorem}

Serre's theorem even writes the algebra down from the diagram. Let $\alpha_1,\dots,\alpha_r$ be the
simple roots and $a_{ij}=n_{\alpha_j\alpha_i}$ the \emph{Cartan matrix}. Then $\mathfrak g$ is generated by
$e_i,f_i,h_i$ ($1\le i\le r$), which are exactly the $\mathfrak{sl}_2$-triples attached to the simple
roots, subject to the relations
\[
  [h_i,h_j]=0,\quad [h_i,e_j]=a_{ij}e_j,\quad [h_i,f_j]=-a_{ij}f_j,\quad [e_i,f_j]=\delta_{ij}h_i,
\]
\[
  (\operatorname{ad}e_i)^{1-a_{ij}}(e_j)=0,\qquad (\operatorname{ad}f_i)^{1-a_{ij}}(f_j)=0\qquad(i\neq j).
\]
The last two relations say that an $\alpha_i$-string through $\alpha_j$ has exactly the length the
Cartan integer predicts. For $\mathfrak{sl}_n$ this recovers $e_i=E_{i,i+1}$, $f_i=E_{i+1,i}$ and
$h_i=E_{ii}-E_{i+1,i+1}$.

The dimension formula $\dim\mathfrak g=\operatorname{rank}+\abs{\Delta}$ and the root counts ($12$ for
$G_2$ from Lecture~4; $48$, $126$, $240$ from Lecture~6; and $72$ for $E_6$, which one checks by counting
the roots of $E_8$ orthogonal to an $A_2$) give the dimensions of the exceptional algebras:
\begin{center}
\renewcommand{\arraystretch}{1.2}
\begin{tabular}{lccccc}
\toprule
 & $G_2$ & $F_4$ & $E_6$ & $E_7$ & $E_8$\\
\midrule
rank & 2 & 4 & 6 & 7 & 8\\
$\abs{\Delta}$ & 12 & 48 & 72 & 126 & 240\\
$\dim\mathfrak g$ & 14 & 52 & 78 & 133 & 248\\
\bottomrule
\end{tabular}
\end{center}
In particular, the $240$ minimal vectors of the $E_8$ lattice, the kissing configuration in
dimension $8$, \emph{literally} label the root spaces of the largest exceptional Lie algebra,
$\mathfrak e_8=\mathfrak h\oplus\bigoplus_{\alpha\in E_8}\mathfrak g_\alpha$, $248=8+240$. The smallest one,
$\mathfrak g_2$, is the algebra of derivations of the octonions (\'E.~Cartan).

%======================================================================
\section{From the algebra to the group: the exponential map}
\label{LIE:sec:exp}

The name ``Lie algebra'' comes from Sophus Lie's idea that a group of transformations is determined,
near the identity, by its \emph{infinitesimal} transformations. For matrices this can be made
completely explicit.

\begin{definition}
For $X\in M_n$ (real or complex) the \emph{matrix exponential} is
\[
  \exp X=e^X=\sum_{k=0}^\infty\frac{X^k}{k!}.
\]
The series converges absolutely, since $\norm{X^k}\le\norm{X}^k$ in the operator norm.
\end{definition}

\begin{proposition}\label{LIE:prop:exp}
For $X,Y\in M_n$ and $g\in\GL_n$:
\begin{enumerate}[nosep]
\item if $XY=YX$, then $e^{X+Y}=e^Xe^Y$; in particular $e^{(s+t)X}=e^{sX}e^{tX}$ and $(e^X)^{-1}=e^{-X}$;
\item $\frac{d}{dt}e^{tX}=Xe^{tX}=e^{tX}X$;
\item $e^{gXg^{-1}}=g\,e^Xg^{-1}$ and $(e^X)^T=e^{X^T}$;
\item $\det e^X=e^{\operatorname{tr}X}$;
\item $\exp$ maps a neighbourhood of $0$ diffeomorphically onto a neighbourhood of $I$, with inverse
  $\log(I+A)=\sum_{k\ge1}(-1)^{k+1}A^k/k$ for $\norm{A}<1$.
\end{enumerate}
\end{proposition}

\begin{proof}[Sketch]
(1) For commuting $X$ and $Y$ the binomial theorem applies to $(X+Y)^k$, and the Cauchy product of the
two series gives the result. (2) Differentiate the series term by term. (3) Apply the identities
$(gXg^{-1})^k=gX^kg^{-1}$ and $(X^k)^T=(X^T)^k$ to each term. (4) Over $\C$, write $X=gTg^{-1}$ with
$T$ upper triangular (Schur). Then $e^T$ is upper triangular with diagonal entries $e^{t_{ii}}$, so
$\det e^X=\det e^T=e^{\sum t_{ii}}=e^{\operatorname{tr}X}$. (5) The derivative of $\exp$ at $0$ is the
identity, so the inverse function theorem applies; the series for $\log$ inverts $\exp$ near $I$ as a
formal identity of power series.
\end{proof}

By (1), $t\mapsto e^{tX}$ is a homomorphism $(\R,+)\to\GL_n$: a \emph{one-parameter subgroup}. This is
the matrix version of the one-parameter groups of diffeomorphisms of Lecture~9. The linear vector field $v\mapsto Xv$ on $\R^n$ generates the
one-parameter group of diffeomorphisms $\varphi_t(v)=e^{tX}v$, because $\frac{d}{dt}e^{tX}v=X(e^{tX}v)$
by (2). So $X$ is the ``velocity'' of the subgroup at the identity.

\begin{figure}[ht]
\centering
\begin{tikzpicture}[scale=1.4]
  \draw[thick] (0,0) circle (1);
  \draw[addcol,thick] (1,-1.35)--(1,1.35) node[above,font=\small]{$T_I\SO(2)=\mathfrak{so}_2\cong\R$};
  \fill (1,0) circle (0.03) node[right,font=\small]{$I$};
  \foreach \t in {0.6,1.2} {\fill[addcol] (1,\t) circle (0.03); \fill ({cos(\t r)},{sin(\t r)}) circle (0.03);
     \draw[-{Stealth},gray,dashed] (0.97,\t) to[bend right=20] ({0.97*cos(\t r)},{0.97*sin(\t r)});}
  \node[font=\small] at (0,0) {$\SO(2)\cong\Sph^1$};
  \node[font=\scriptsize,gray] at (2.45,0.45) {$\theta\mapsto\begin{pmatrix}\cos\theta&-\sin\theta\\ \sin\theta&\cos\theta\end{pmatrix}$};
\end{tikzpicture}
\caption{The exponential map wraps the tangent line $\mathfrak{so}_2$ around the circle $\SO(2)$.
Its kernel is $2\pi\Z$, a lattice.}
\label{LIE:fig:so2}
\end{figure}

\subsection*{Matrix groups and their Lie algebras}

\begin{definition}
A \emph{matrix group} is a closed subgroup $G\subseteq\GL_n(\R)$ or $\GL_n(\C)$. Its \emph{Lie algebra} is
\[
  \operatorname{Lie}(G)=\{X\in M_n: e^{tX}\in G\text{ for all }t\in\R\}.
\]
\end{definition}

\begin{theorem}[von Neumann 1929; Cartan 1930]\label{LIE:thm:closed}
Let $G$ be a matrix group. Then $\operatorname{Lie}(G)$ is a real vector space closed under the commutator,
so it is a Lie algebra. Moreover $G$ is a smooth submanifold of $M_n$ (a \emph{Lie group}) with tangent
space $T_IG=\operatorname{Lie}(G)$, and $\exp$ maps a neighbourhood of $0$ in $\operatorname{Lie}(G)$
diffeomorphically onto a neighbourhood of $I$ in $G$.
\end{theorem}

We will not prove that $G$ is a manifold (see Hall, Ch.~3, or Stillwell, \emph{Naive Lie Theory}, Ch.~7), but it is worth seeing where the bracket comes from. Closure under addition follows from the
\emph{Lie product formula} $e^{X+Y}=\lim_{k\to\infty}\big(e^{X/k}e^{Y/k}\big)^k$. For the bracket, take
$X,Y\in\operatorname{Lie}(G)$. For each $t$ the matrix $e^{tX}Ye^{-tX}$ lies in $\operatorname{Lie}(G)$, because
$\exp\!\big(s\,e^{tX}Ye^{-tX}\big)=e^{tX}e^{sY}e^{-tX}\in G$. Since $\operatorname{Lie}(G)$ is a closed
subspace, it also contains the derivative
\[
  \frac{d}{dt}\Big|_{t=0}e^{tX}Ye^{-tX}=XY-YX=[X,Y].
\]
Equivalently, expanding the group commutator gives
\[
  e^{sX}e^{tY}e^{-sX}e^{-tY}=I+st\,[X,Y]+(\text{terms of total degree}\ge3\text{ in }s,t).
\]
\emph{The bracket is the infinitesimal shadow of the failure of the group to commute.}

\begin{example}[The classical groups recover the classical algebras]\leavevmode
\begin{itemize}[nosep]
\item $\operatorname{Lie}(\GL_n)=M_n=\mathfrak{gl}_n$.
\item $\operatorname{Lie}(\SL_n)=\mathfrak{sl}_n$. If $\operatorname{tr}X=0$ then $\det e^{tX}=1$ by
  Proposition~\ref{LIE:prop:exp}(4); conversely, $\det e^{tX}=e^{t\operatorname{tr}X}=1$ for all $t$ forces
  $\operatorname{tr}X=0$.
\item $\operatorname{Lie}(G_J)=\mathfrak g_J$ for the isometry group $G_J=\{g: g^TJg=J\}$ of the form $J$. If
  $X^TJ+JX=0$, then $X^T=-JXJ^{-1}$, so $(e^{X})^TJe^{X}=Je^{-X}J^{-1}Je^{X}=J$; applying this to $tX$ shows $\mathfrak g_J\subseteq
  \operatorname{Lie}(G_J)$. Conversely, differentiating $(e^{tX})^TJe^{tX}=J$ at $t=0$ gives $X^TJ+JX=0$. Thus $\mathfrak{so}_N$
  belongs to $\Orth(N)$ and $\SO(N)$, and $\mathfrak{sp}_{2n}$ to $\mathrm{Sp}_{2n}$. This is why $\mathfrak g_J$
  deserves the name ``infinitesimal isometries''.
\item The unitary group $\mathrm U(n)=\{g\in\GL_n(\C): g^*g=I\}$ has Lie algebra
  $\mathfrak u(n)=\{X: X^*=-X\}$ (skew-Hermitian matrices), and $\mathrm{SU}(n)$ has
  $\mathfrak{su}(n)=\mathfrak u(n)\cap\mathfrak{sl}_n$. Note that $\mathfrak{su}(n)\otimes\C=\mathfrak{sl}_n(\C)$: the
  compact group $\mathrm{SU}(n)$ is a \emph{real form} of $\mathfrak{sl}_n(\C)$.
\end{itemize}
\end{example}

\begin{example}[$\SO(2)$ and $\SO(3)$]\label{LIE:ex:so3}
For $X=\theta\left(\begin{smallmatrix}0&-1\\1&0\end{smallmatrix}\right)$ we have $X^2=-\theta^2I$, and
splitting the series into even and odd terms gives
\[
  e^X=\cos\theta\,I+\sin\theta\begin{pmatrix}0&-1\\1&0\end{pmatrix},
\]
which is rotation by $\theta$ (Figure~\ref{LIE:fig:so2}). In $\SO(3)$, write $X_\omega v=\omega\times v$
with $\theta=\abs{\omega}$. Using $X_\omega^3=-\theta^2X_\omega$, the same splitting gives
\emph{Rodrigues' formula}
\[
  e^{X_\omega}=I+\frac{\sin\theta}{\theta}X_\omega+\frac{1-\cos\theta}{\theta^2}X_\omega^2,
\]
which is rotation by angle $\theta$ about the axis $\omega$. Every rotation has an axis, so $\exp$ is
onto $\SO(3)$. It maps the closed ball $\abs{\omega}\le\pi$ onto $\SO(3)$, identifying only antipodal
points of the boundary sphere, so $\SO(3)\cong\R P^3$.
\end{example}

\begin{remark}[Same algebra, different groups]
The group $\mathrm{SU}(2)$ of unit quaternions is a $3$-sphere. It acts on the purely imaginary
quaternions $\cong\R^3$ by conjugation, and this gives a two-to-one homomorphism
$\mathrm{SU}(2)\to\SO(3)$ with kernel $\{\pm I\}$. The two groups have isomorphic Lie algebras,
$\mathfrak{su}(2)\cong\mathfrak{so}(3)\cong(\R^3,\times)$, but they are not isomorphic: one is simply connected,
the other has $\pi_1=\Z/2$. The Lie algebra determines a connected group only up to coverings. The
missing global information turns out to be a \emph{lattice}, as the next section shows.
\end{remark}

%======================================================================
\section{Tori, Weyl groups, and the lattice inside the group}
\label{LIE:sec:torus}

Return to $A_{n-1}$, now in the compact group $\mathrm{SU}(n)$. The diagonal matrices in $\mathrm{SU}(n)$ form
the \emph{maximal torus}
\[
  T=\big\{\operatorname{diag}(e^{i\theta_1},\dots,e^{i\theta_n}):\ \theta_1+\dots+\theta_n\equiv0\pmod{2\pi}\big\},
\]
with Lie algebra
\[
  \mathfrak t=\{\operatorname{diag}(i\theta_1,\dots,i\theta_n):\ \theta_1+\dots+\theta_n=0\}
  \cong\{\theta\in\R^n:\textstyle\sum\theta_j=0\},
\]
which is again the hyperplane of Lecture~6. Since $\mathfrak t$ is commutative, $\exp\colon\mathfrak t\to T$ is a
surjective homomorphism, and its kernel consists of the $\theta$ with every $\theta_j\in2\pi\Z$ and
$\sum\theta_j=0$:
\[
  \ker\big(\exp|_{\mathfrak t}\big)=2\pi\cdot\{x\in\Z^n:\textstyle\sum x_j=0\}=2\pi\cdot A_{n-1},
\]
the $A_{n-1}$ root lattice from Lecture~6. So
\[
  T\;\cong\;\mathfrak t\,/\,2\pi A_{n-1}:
\]
\emph{the maximal torus of $\mathrm{SU}(n)$ is Euclidean space modulo the $A_{n-1}$ root lattice.} For
$\mathrm{SU}(3)$ it is the plane modulo the hexagonal lattice, the lattice of the densest disk packing
from Lecture~1.

\paragraph{The Weyl group.} Permutation matrices (with one sign changed if necessary, to make the
determinant $1$) normalize $T$ and permute the $\theta_j$. The resulting group
\[
  W=N(T)/T\cong\Sigma_n
\]
is the Weyl group of $A_{n-1}$, generated by the reflections $s_{e_i-e_j}$, which swap $\theta_i$ and
$\theta_j$. By the spectral theorem every element of $\mathrm{SU}(n)$ is conjugate to an element of $T$,
uniquely up to the action of $W$, so
\[
  \{\text{conjugacy classes of }\mathrm{SU}(n)\}\;\longleftrightarrow\;T/W
  \;\longleftrightarrow\;\{\text{unordered $n$-tuples of eigenvalues on }\Sph^1\text{ with product }1\}.
\]

\begin{remark}[Back to configuration spaces]
An element of $\mathrm{SU}(n)$ is \emph{regular} when its eigenvalues are distinct, that is, when its
$n$ eigenvalues form a configuration of distinct points on the circle. Choose real eigenvalue angles with
$\sum\theta_j=0$, ordered so that $\theta_1\ge\theta_2\ge\dots\ge\theta_n\ge\theta_1-2\pi$ (this is
always possible, in exactly one way), and record the gaps
\[
  g_k=\theta_k-\theta_{k+1}\ (1\le k<n),\qquad g_n=2\pi-(\theta_1-\theta_n).
\]
The gaps are non-negative and sum to $2\pi$, and together with $\sum\theta_j=0$ they determine the
$\theta_j$. So $T/W$ is a closed simplex of dimension $n-1$, the \emph{Weyl alcove}. Its interior, the
regular classes, is identified through the gaps with the reduced configuration space of $n$ points on a
circle in a fixed cyclic order (one of the open simplices of gaps in Lecture~7). For $n=3$ it is the same open triangle of gap coordinates $(x,y,1-x-y)$, up to the scale
factor $2\pi$. Its boundary walls, where two eigenvalues collide, are the walls $\theta_i=\theta_j$ of
$\Conf(n,\R)$ from Lecture~8, wrapped around the circle.
\end{remark}

All of this generalizes.

\begin{theorem}[Maximal tori; Weyl, Cartan]\label{LIE:thm:torus}
Let $G$ be a compact connected Lie group. Then $G$ has a maximal torus $T$, unique up to conjugacy,
and every element of $G$ is conjugate into $T$. The Weyl group $W=N(T)/T$ is finite and is the Weyl
group of the root system of $\mathfrak g\otimes\C$. The kernel of $\exp\colon\mathfrak t\to T$ is a lattice
containing $2\pi$ times the coroot lattice (the coroots $h_\alpha$ placed in $\mathfrak t$ as $ih_\alpha$,
as in the $\mathrm{SU}(n)$ example), with equality exactly when $G$ is simply connected; in general the
quotient is $\pi_1(G)$. The
compact simply connected simple Lie groups correspond to the connected Dynkin diagrams:
\[
  \mathrm{SU}(n+1)\ (A_n),\quad \mathrm{Spin}(2n+1)\ (B_n),\quad \mathrm{Sp}(n)\ (C_n),\quad
  \mathrm{Spin}(2n)\ (D_n),\quad G_2,\ F_4,\ E_6,\ E_7,\ E_8 .
\]
\end{theorem}

For example, $\mathrm{SU}(2)$ has $T=\{\operatorname{diag}(e^{i\theta},e^{-i\theta})\}$ and kernel
$\theta\in2\pi\Z$, while $\SO(3)=\mathrm{SU}(2)/\{\pm I\}$ has the finer kernel $\theta\in\pi\Z$. The covering
$\mathrm{SU}(2)\to\SO(3)$ of the previous section is recorded by an index-two inclusion of lattices.

%======================================================================
\section{Summary: a dictionary}
\label{LIE:sec:dictionary}

\begin{center}
\renewcommand{\arraystretch}{1.25}
\begin{tabular}{p{0.42\textwidth}p{0.50\textwidth}}
\toprule
Packings and lattices (Lectures 4--8) & Lie theory\\
\midrule
root system $\Delta$ & nonzero eigenvalues of $\operatorname{ad}\mathfrak h$ on $\mathfrak g$\\
Cartan integer $n_{\beta\alpha}\in\Z$ & eigenvalue of the coroot $h_\alpha$; integral by $\mathfrak{sl}_2$ theory\\
reflection $s_\alpha$, Weyl group $W$ & symmetry of $\alpha$-strings; $W=N(T)/T$\\
choice of $d$, positive roots, Weyl chamber & upper triangular matrices (a Borel subalgebra)\\
simple roots, Dynkin diagram & generators $e_i,f_i,h_i$ and Serre relations\\
rank $+\abs{\Delta}$ & $\dim\mathfrak g$\\
root lattice $\Z\Delta$ (coroot lattice, in general) & $\tfrac1{2\pi}\,$kernel of $\exp$ on the maximal torus (simply connected, compact)\\
$E_8$ kissing configuration ($240$ vectors) & root spaces of $\mathfrak e_8$, $248=8+240$\\
walls $x_i=x_j$ of $\Conf(N,\R)$ & eigenvalue collisions; non-regular elements of $\mathrm{SU}(N)$\\
one-parameter groups of diffeomorphisms (Lecture 9) & $t\mapsto e^{tX}$\\
\bottomrule
\end{tabular}
\end{center}

%======================================================================
\section*{Exercises}
\addcontentsline{toc}{section}{Exercises}

\begin{exercise}
Write out the root decomposition of $\mathfrak{sl}_3$ explicitly. Check that $[E_{12},E_{23}]=E_{13}$,
$[E_{13},E_{32}]=E_{12}$ and $[E_{12},E_{13}]=0$, and interpret each identity on the hexagon of
Example~\ref{LIE:ex:sl3}.
\end{exercise}

\begin{exercise}
Prove the Jacobi identity for the commutator, and show that $\omega\mapsto X_\omega$ is an isomorphism
$(\R^3,\times)\to\mathfrak{so}_3(\R)$.
\end{exercise}

\begin{exercise}
Expand $e^{sX}e^{tY}e^{-sX}e^{-tY}$ to second order and confirm that the result is $I+st[X,Y]$ plus
higher-order terms.
\end{exercise}

\begin{exercise}
Show that $\exp\colon\mathfrak{sl}_2(\R)\to\SL_2(\R)$ is not surjective: the matrix
$\left(\begin{smallmatrix}-1&1\\0&-1\end{smallmatrix}\right)$ is not in the image. \emph{Hint:} a real
traceless $2\times2$ matrix has eigenvalues $\pm\lambda$ or $\pm i\mu$. Compute the possible traces of
$e^X$, and decide when the trace can equal $-2$.
\end{exercise}

\begin{exercise}
Using the antidiagonal form $J$, find the roots of $\mathfrak{so}_6$ and check that they form a copy of
$A_3$ (twelve roots, simply laced, rank $3$). Do the same for $\mathfrak{so}_4$ and $A_1\times A_1$, and for
$\mathfrak{sp}_4$ and $\mathfrak{so}_5$ (show that $C_2$ and $B_2$ differ by a rotation of $45^\circ$ and a scaling).
\end{exercise}

\begin{exercise}
Prove Rodrigues' formula, and show that $\exp$ maps the closed ball of radius $\pi$ in $\mathfrak{so}_3$
onto $\SO(3)$, identifying exactly the antipodal points of the boundary.
\end{exercise}

\begin{exercise}
Draw the kernel of $\exp$ on the maximal torus of $\mathrm{SU}(3)$ inside the plane
$\theta_1+\theta_2+\theta_3=0$, together with the Weyl alcove. Check that the kernel is a hexagonal
lattice and that the alcove is an equilateral triangle. How many alcoves fit into a fundamental domain
of the lattice?
\end{exercise}

\section*{Further reading}
\addcontentsline{toc}{section}{Further reading}
\begin{itemize}[nosep]
\item J.~Stillwell, \emph{Naive Lie Theory}, Springer UTM, 2008. Matrix groups first, with minimal prerequisites.
\item B.~C.~Hall, \emph{Lie Groups, Lie Algebras, and Representations}, 2nd ed., Springer GTM 222, 2015.
  Matrix groups, the exponential map, and a full treatment of the classification.
\item J.~E.~Humphreys, \emph{Introduction to Lie Algebras and Representation Theory}, Springer GTM 9, 1972.
  The algebraic route from $\mathfrak{sl}_2$ to root systems and Serre's theorem.
\item W.~Fulton and J.~Harris, \emph{Representation Theory: A First Course}, Springer GTM 129, 1991.
  Lectures 11--13 work through $\mathfrak{sl}_2$ and $\mathfrak{sl}_3$ by hand, and Lecture 21 does the classification.
\item J.~H.~Conway and N.~J.~A.~Sloane, \emph{Sphere Packings, Lattices and Groups}, Ch.~4 and 8, for the
  lattice side of $E_8$.
\end{itemize}
